% Sampul belakang, versi formal2. Dipanggil dari main.tex lewat \versiSampul.
% Latar kertas, pita seri di atas, blok sinopsis, daftar topik, dan panel marun di bawah dengan kotak-kotak "kode".
% Pasangan di sampul depan: depan-formal2.tex. Tanpa logo.
\clearpage
\thispagestyle{empty}
\begin{tikzpicture}[remember picture,overlay]
  % --- Latar kertas ---
  \fill[kertas] (current page.south west) rectangle (current page.north east);

  % --- Pita seri di atas (sama dengan sampul depan) ---
  \fill[marun] (current page.north west) rectangle ([yshift=-16mm]current page.north east);
  \fill[emas] ([yshift=-16mm]current page.north west) rectangle ([yshift=-17.2mm]current page.north east);
  \node[anchor=west,text=emasmuda,font=\capsfont\small,inner sep=0pt]
    at ([xshift=20mm,yshift=-8mm]current page.north west) {\jenisBuku};
  \node[anchor=east,text=emas,font=\judul\fontsize{22}{22}\selectfont,inner sep=0pt]
    at ([xshift=-20mm,yshift=-8.5mm]current page.north east) {\{\,\}};

  % --- Judul ringkas ---
  \fill[emas] ([xshift=20mm,yshift=-34mm]current page.north west) rectangle ++(14mm,-1.2pt);
  \node[anchor=west,text=merah,font=\capsfont\small,inner sep=0pt]
    at ([xshift=38mm,yshift=-34.2mm]current page.north west) {Tentang buku};
  \node[anchor=north west,text=gelap,font=\judul\fontsize{34}{38}\selectfont,align=left,inner sep=0pt,
        text width=165mm]
    (judul) at ([xshift=20mm,yshift=-46mm]current page.north west) {\hyphenpenalty=10000 \exhyphenpenalty=10000 \judulBuku\par};
  \node[anchor=north west,text=merah,font=\itshape\fontsize{15}{19}\selectfont,align=left,inner sep=0pt,
        text width=150mm]
    (sub) at ([yshift=-5mm]judul.south west) {\subJudul};

  % --- Sinopsis ---
  \node[anchor=north west,text=gelap,font=\sffamily\fontsize{11.5}{18}\selectfont,align=left,inner sep=0pt,
        text width=160mm]
    (isi) at ([yshift=-12mm]sub.south west)
    {Buku ini membahas pemrograman berorientasi objek dengan Python secara bertahap, mulai dari pengenalan paradigma,
     class dan object, hingga prinsip desain dan penerapannya dalam aplikasi.\par\vspace{8pt}
     Setiap bab dilengkapi penjelasan konsep, contoh kode beserta hasil eksekusinya, studi kasus, dan latihan,
     serta ditutup dengan pengembangan proyek aplikasi yang memadukan seluruh konsep yang telah dipelajari.};

  % --- Topik utama ---
  \node[anchor=north west,text=merah,font=\capsfont\small,inner sep=0pt]
    (topik) at ([yshift=-14mm]isi.south west) {Topik utama};
  \fill[emas] ([yshift=-2.5mm]topik.south west) rectangle ++(14mm,-1.2pt);
  \node[anchor=north west,text=gelap,font=\sffamily\fontsize{11.5}{19}\selectfont,align=left,inner sep=0pt]
    at ([yshift=-8mm]topik.south west)
    {%
      \begin{minipage}[t]{78mm}\raggedright
        {\color{merah}\textbullet}\ Class dan object\\
        {\color{merah}\textbullet}\ Encapsulation dan decorator\\
        {\color{merah}\textbullet}\ Inheritance\\
        {\color{merah}\textbullet}\ Abstraksi dan polymorphism\\
        {\color{merah}\textbullet}\ Relasi antarkelas
      \end{minipage}\hfill
      \begin{minipage}[t]{78mm}\raggedright
        {\color{merah}\textbullet}\ Modular programming\\
        {\color{merah}\textbullet}\ Exception handling\\
        {\color{merah}\textbullet}\ File handling\\
        {\color{merah}\textbullet}\ Antarmuka grafis (GUI)\\
        {\color{merah}\textbullet}\ Proyek aplikasi terpadu
      \end{minipage}%
    };

  % --- Panel marun di bawah ---
  \fill[marun] (current page.south west) rectangle ([yshift=78mm]current page.south east);
  \fill[emas] ([yshift=78mm]current page.south west) rectangle ([yshift=79.2mm]current page.south east);

  % Kotak-kotak kode di sisi kanan panel
  \pgfmathsetseed{11}
  \foreach \i in {1,...,110}{
    \pgfmathsetmacro{\px}{200-100*rnd*rnd}
    \pgfmathsetmacro{\py}{6+66*rnd}
    \pgfmathsetmacro{\sz}{0.8+rnd*rnd*6}
    \pgfmathsetmacro{\op}{0.10+0.55*rnd}
    \fill[emasmuda,opacity=\op]
      ([xshift=\px mm,yshift=\py mm]current page.south west) rectangle ++(\sz mm,\sz mm);
  }
  \foreach \i in {1,...,24}{
    \pgfmathsetmacro{\px}{200-100*rnd*rnd}
    \pgfmathsetmacro{\py}{6+66*rnd}
    \pgfmathsetmacro{\sz}{0.6+rnd*2.2}
    \fill[merah,opacity=0.9]
      ([xshift=\px mm,yshift=\py mm]current page.south west) rectangle ++(\sz mm,\sz mm);
  }

  % Instansi dan tahun
  \node[anchor=south west,text=emasmuda,font=\capsfont\scriptsize,align=left,inner sep=0pt,text width=120mm]
    at ([xshift=20mm,yshift=12mm]current page.south west)
    {\prodi\\[3pt]\fakultas\\[3pt]\universitas\\[3pt]\tahun};
\end{tikzpicture}
\clearpage